%% AISSH-26 manuscript, built on the official sn-jnl.cls template
%% (paper/Archive/sn-jnl.cls, matching the sn-mathphys-num option properly)
%% instead of the older GitHub-mirror class file used by the other drafts in
%% this folder. Content matches aissh_springer_manuscript_v2.tex (see that
%% file's header for full narrative-restructure rationale): guard-bear's
%% safety-transfer failure leads, consumer-hardware deployability second, and
%% the LoRA configuration-ordering crossover is framed as an unexpected
%% discovery made during development rather than the primary motivation.
%% Body condensed further so body+references together fit 12 pages total
%% (the venue's stated limit, references included). Bibliography uses BibTeX
%% (aissh_final.bib) instead of a manually-numbered list.

\documentclass[pdflatex,sn-mathphys-num]{sn-jnl-official}% official class, numbered Math/Phys reference style

\usepackage{graphicx}%
\usepackage{multirow}%
\usepackage{amsmath,amssymb,amsfonts}%
\usepackage{amsthm}%
\usepackage{mathrsfs}%
\usepackage{xcolor}%
\usepackage{textcomp}%
\usepackage{manyfoot}%
\usepackage{booktabs}%
\usepackage{enumitem}%

\raggedbottom
\renewcommand{\arraystretch}{0.85}
\setlist{nosep}

\begin{document}

\title[Safety-Gated, Age-Adaptive LoRA Fine-Tuning for a Virtual Pediatric Hospital Guide]{Safety-Gated, Age-Adaptive LoRA Fine-Tuning for a Virtual Pediatric Hospital Guide}

\author*[1]{\fnm{Anthony} \sur{Barros}}\email{abarros6@uwo.ca}

\author[1]{\fnm{Kelsey} \sur{Kloosterman}}\email{kkloost2@uwo.ca}

\author[2,3]{\fnm{Sandrine} \sur{de Ribaupierre}}\email{sderibau@uwo.ca}

\author[1,2]{\fnm{Roy} \sur{Eagleson}}\email{eagleson@uwo.ca}

\affil*[1]{\orgdiv{Department of Artificial Intelligence and Software Engineering, Faculty of Engineering}, \orgname{Western University}, \orgaddress{\street{1151 Richmond Street}, \city{London}, \postcode{N6A3K7}, \state{Ontario}, \country{Canada}}}

\affil[2]{\orgdiv{Centre for Brain and Mind}, \orgname{Western University}, \orgaddress{\street{1151 Richmond Street}, \city{London}, \postcode{N6A3K7}, \state{Ontario}, \country{Canada}}}

\affil[3]{\orgdiv{Clinical Neurological Sciences}, \orgname{Western University}, \orgaddress{\street{1151 Richmond Street}, \city{London}, \postcode{N6A3K7}, \state{Ontario}, \country{Canada}}}

\abstract{Pediatric patients navigating hospital environments need age-appropriate, institutionally grounded information that general-purpose AI systems are not designed to provide, and any system accepting free-text input from children must defend against inputs it was never designed to handle. We present \textit{small-bear}, an age-targeted conversational guide deployed within a digital recreation of Victoria Hospital (London, Ontario), paired with our \textit{guard-bear}, an upstream input-safety classifier -- together a complete pipeline running entirely on consumer hardware (an Apple Mac Mini M4, 16~GB) with no cloud dependency. The first finding concerns safety: the unmodified base classifier we build on, \texttt{Prompt-Guard-86M}~\cite{promptguard2024}, performs \textit{worse than random} (ROC-AUC 0.466) on pediatric-context queries, saturating by flagging nearly everything as unsafe. Fine-tuning \textit{guard-bear} on 4{,}703 pediatric-specific labeled examples resolves this to near-perfect discrimination (ROC-AUC 0.999, 0.987 recall, 0.990 precision) -- evidence that safety mechanisms validated for a developer-facing threat model do not transfer to a different population without dedicated evaluation. The second finding concerns \textit{small-bear}: we fine-tune Llama 3.2 Instruct (1B, 3B) with LoRA at two configurations, Standard ($r=8$, 16 layers) and Fast ($r=4$, 8 layers), on 1123 LIMA-inspired~\cite{zhou2023lima} synthetic examples, substantially outperforming zero-shot base models on readability (56--66\% FK~$\leq$~7.0 pass rate vs.\ 12\%) and inter-role lexical separation (0.92--0.96 vs.\ 0.64--0.67), with Fast-1B meeting the 1.0-second real-time VR latency target. During development, an unexpected \textbf{configuration-ordering crossover} emerged (Standard better on 3B, Fast better on 1B); since run-to-run variance can exceed that gap, we validated it with a dedicated 110-run multi-seed campaign, confirming it is statistically real on both sides, though length-adjustment shows the 3B-side gap is largely a length effect while the 1B-side gap is length-independent. Together, these results show a deployable, age-adaptive pediatric pipeline is achievable on consumer hardware, and its safety layer, not response quality, posed the larger practical risk when unvalidated.}

\keywords{parameter-efficient fine-tuning, low-rank adaptation, pediatric clinical NLP, virtual reality, on-device inference, small language models, LLM input safety classification}

\maketitle

\section{Introduction}\label{sec1}

Large language models (LLMs) enable domain-specific conversational systems, but resource-constrained clinical deployment raises constraints rarely addressed together: limited memory favors smaller models, domain adaptation requires fine-tuning on specialized data, real-time interaction imposes strict latency bounds, and free-text input from a vulnerable population must be defended against inputs the system was never designed to handle.

This paper presents \textit{small-bear}, an age-targeted conversational guide embedded in a virtual recreation of Victoria Hospital (London, Ontario), for children aged 5--11 and adolescents aged 12--18 within a VR environment, together with \textit{guard-bear}, an upstream input-safety classifier -- both named after LHSC's real-world bear mascot. All training and inference runs on an Apple Mac Mini M4 (16~GB): a complete pipeline deployable on consumer hardware without cloud infrastructure.

A conversational system that accepts open-ended text from children needs a defense layer distinct from \textit{small-bear} itself. General-purpose prompt-injection classifiers such as \texttt{Prompt-Guard-86M}~\cite{promptguard2024} target a developer-facing threat model that does not describe a pediatric system's actual inputs: short, first-person legitimate queries on one side, gibberish and adult content on the other. Our \textit{guard-bear} closes this gap, finding, as our central safety result, that the unmodified base classifier performs worse than random here, a failure fine-tuning fully resolves.

For \textit{small-bear}, we fine-tune Llama 3.2 Instruct models~\cite{dubey2024llama3} at 1B and 3B scales using LoRA~\cite{hu2022lora}, alongside QLoRA~\cite{dettmers2023qlora}. Training data follows the LIMA hypothesis~\cite{zhou2023lima}: $\sim$560 examples per adapter (1123 total) from Victoria Hospital's public materials. We compare Standard ($r=8$, 16 layers) and Fast ($r=4$, 8 layers) adapters across both sizes and age groups. Pediatric deployment introduces demands distinct from adult-facing clinical LLMs and passive VR tools -- developmental vocabulary, supportive register, hard safety constraints; we are not aware of prior work fine-tuning age-stratified LoRA adapters for this setting (Section~\ref{sec:relatedwork}).

During development, an initial single-seed ablation suggested an unexpected \textbf{configuration-ordering crossover}: Standard yields better readability on 3B, Fast on 1B -- opposite what existing rank-selection guidance, developed only at the 7B--70B scale, predicts. A reproducibility check found run-to-run variance exceeding the crossover gap, so we validated it with a dedicated 110-run multi-seed campaign (Section~\ref{sec:seedcampaign}). A subsequent rank sweep (Section~\ref{sec:ranksweep}) attributed it to rank, replicating cleanly for SmolLM2 but only weakly for others.

\medskip\noindent The contributions of this paper are:
\begin{itemize}
    \item \textit{guard-bear}: a fine-tuned input-safety classifier showing its unmodified general-purpose base model performs \textit{worse than random} on this population -- resolved by fine-tuning to near-perfect discrimination -- evidence that safety classifiers cannot be assumed to transfer to a new population without dedicated evaluation.
    \item A complete age-adaptive pipeline, response generation gated by an upstream safety classifier, trained and evaluated end-to-end on consumer hardware with no cloud dependency.
    \item An unexpected LoRA configuration-ordering crossover between Llama 3.2 1B and 3B, discovered during development and confirmed statistically real on both sides by a dedicated 110-run multi-seed campaign, after an initial single-seed observation proved non-reproducible at face value (Section~\ref{sec:seedcampaign}).
    \item A rank sweep and layer sweep investigating the crossover's mechanism (Section~\ref{sec:ranksweep}), and cross-architecture evidence (Qwen~2, SmolLM2, Qwen 2.5) that the Fast-vs-Standard directional pattern generalizes beyond Llama even where the rank-regime mechanism does not cleanly resolve.
\end{itemize}

\section{Related Work}\label{sec:relatedwork}

General-purpose prompt-injection classifiers such as \texttt{Prompt-Guard-86M}~\cite{promptguard2024} target developer-facing override patterns; the closest child-safety effort, \textit{KidRails}~\cite{kidrails2025}, fully fine-tunes an 8B response model without a separate classifier or baseline comparison. \textit{guard-bear} instead evaluates a dedicated input classifier against its base model on a labeled pediatric threat taxonomy -- to our knowledge, the first such evaluation for this population.

Prompted, not fine-tuned, LLMs have been evaluated as pediatric communication aids: Rao et al.~\cite{rao2025pediatric} found off-the-shelf models can produce age-appropriate surgical explanations but called for further fine-tuning; a Norwegian evaluation similarly found generated child-directed language skewed too advanced~\cite{hassan2025childconv}. \textit{KidLM}~\cite{nayeem2024kidlm} adapts masked-language-model pretraining to a children's corpus but targets general modeling, not clinical generation. VR-based interventions reduce procedural anxiety~\cite{ryu2021vr} but are largely passive, scripted content. No prior work, to our knowledge, fine-tunes age-stratified LoRA adapters for pediatric hospital communication.

LoRA~\cite{hu2022lora} and QLoRA~\cite{dettmers2023qlora} are standard for resource-constrained fine-tuning, but rank-selection guidance -- including Biderman et al.'s~\cite{biderman2024lora} empirical study, the most thorough to date -- is validated only at the 7B--70B scale; rank-ordering below 7B is uncharacterized by prior work. Our rank and layer sweeps (Section~\ref{sec:ranksweep}) address this gap.

\section{Results}\label{sec2}

This section first establishes \textit{small-bear} and \textit{guard-bear} as deployable components (Sections~\ref{sec:responsemodel}--\ref{sec:deployment}), then presents the statistical validation and mechanistic investigation of an unexpected LoRA effect discovered during development (Sections~\ref{sec:seedcampaign}--\ref{sec:crossarch}). Methods and dataset detail: Section~\ref{sec:methodology}.

\subsection{small-bear: Perplexity, Readability, Latency, and Style Separation}\label{sec:responsemodel}

Table~\ref{tab:res2} reports a representative single run (seed 42) of the four fine-tuned configurations and both base models. All four fine-tuned adapters substantially outperform base models on FK~$\leq$~7.0 for the 5--11 group (56--66\% vs.\ 12\%) and inter-role style separation (0.92--0.96 vs.\ 0.64--0.67) -- though this classifier does not discriminate Standard from Fast, consistent with, not independent confirmation of, the readability pattern below. \textbf{Fast-1B is the only Llama configuration meeting the 1.0s VR latency target} (0.97s avg); all 3B configurations exceed 2.2s. Post-hoc perplexity does not track FK-based readability cleanly: Fast has lower perplexity than Standard at \emph{both} sizes.

On this seed, Fast-1B beats Standard-1B on FK pass rate (62\% vs.\ 56\%), while Standard-3B and Fast-3B tie (66\% vs.\ 66\%); Section~\ref{sec:seedcampaign} shows this single-run picture undersells the true, smaller 3B-side gap.

\begin{table*}[!h]
  \caption{Readability (50 age-matched held-out prompts/group) and latency/throughput (sequential, Mac Mini M4), single run (seed 42), current dataset vintage.}
  \label{tab:res2}
  \scriptsize
  \setlength{\tabcolsep}{3pt}
  \begin{tabular}{@{}llcccccccc@{}}
    \toprule
    Config. & Age & FK & SMOG & G.Fog & C-L & TTR & FK$\leq$7.0 & Avg Lat.(s) & $<$1.0s \\
    \midrule
    Standard-1B & 5--11  & 7.1  & 9.2  & 8.7  & 7.1  & 0.756 & 56\% & 1.12 & 26\% \\
    Standard-1B & 12--18 & 10.3 & 12.7 & 13.0 & 10.7 & 0.734 & ---  & 1.36 & 2\%  \\
    Standard-3B & 5--11  & 6.6  & 8.9  & 8.1  & 7.1  & 0.781 & 66\% & 2.22 & 0\%  \\
    Standard-3B & 12--18 & 10.4 & 12.8 & 12.9 & 10.6 & 0.733 & ---  & 3.08 & 0\%  \\
    Fast-1B     & 5--11  & 6.5  & 9.0  & 8.0  & 6.8  & 0.752 & 62\% & 0.97 & 62\% \\
    Fast-1B     & 12--18 & 9.7  & 12.2 & 12.1 & 10.0 & 0.725 & ---  & 1.11 & 38\% \\
    Fast-3B     & 5--11  & 6.5  & 9.0  & 8.1  & 7.1  & 0.724 & 66\% & 2.70 & 0\%  \\
    Fast-3B     & 12--18 & 10.0 & 12.5 & 12.5 & 10.4 & 0.663 & ---  & 3.63 & 0\%  \\
    Base-1B     & 5--11  & 9.5  & 11.9 & 12.0 & 9.6  & 0.582 & 12\% & 1.90 & 16\% \\
    Base-1B     & 12--18 & 10.4 & 12.9 & 13.1 & 11.6 & 0.553 & ---  & 2.23 & 2\%  \\
    Base-3B     & 5--11  & 9.8  & 12.4 & 12.5 & 10.2 & 0.609 & 12\% & 4.54 & 8\%  \\
    Base-3B     & 12--18 & 11.7 & 13.6 & 14.7 & 11.4 & 0.551 & ---  & 5.40 & 2\%  \\
    \bottomrule
  \end{tabular}
\end{table*}

\subsection{guard-bear: Safety Classification Results}\label{sec:guardresults}

The clearest evidence that general-purpose safety tooling cannot be assumed to transfer to a new population comes from comparing \textit{guard-bear} against unmodified \texttt{Prompt-Guard-86M} (default threshold 0.50) on the 613-example held-out test set (Table~\ref{tab:guardheadline}). \textbf{The base model is worse than random}: ROC-AUC 0.466 is below chance, since its imperative-override feature space does not describe short, first-person pediatric queries or gibberish/out-of-scope/adult content. Its apparent 0.9585 recall is an artifact of saturation, flagging 294 of 300 safe examples as unsafe too (precision 0.505); a sub-0.5 AUC means the base model ranks true-unsafe examples below true-safe examples across the full threshold range. \textbf{Fine-tuning resolves the discrimination failure}: ROC-AUC rises to 0.999, false negatives fall from 13 to 4, false positives from 294 to 3, meeting all four deployment targets. The largest subcategory gain is \texttt{gibberish\_malformed} (48\%~$\to$~100\% recall).

\begin{table}[!h]
  \centering
  \caption{guard-bear vs.\ base model, held-out test set (613 examples: 300 safe, 313 unsafe).}
  \label{tab:guardheadline}
  \small
  \begin{tabular}{@{}lccc@{}}
    \toprule
    Metric & Base (t=0.50) & guard-bear (t=0.08) & $\Delta$ \\
    \midrule
    Unsafe recall    & 0.9585 & \textbf{0.9872} & $+$0.029 \\
    Unsafe precision & 0.5051 & \textbf{0.9904} & $+$0.485 \\
    Unsafe F1        & 0.6615 & \textbf{0.9888} & $+$0.327 \\
    Accuracy         & 0.4992 & \textbf{0.9886} & $+$0.489 \\
    ROC-AUC          & 0.4660 & \textbf{0.9993} & $+$0.533 \\
    \bottomrule
  \end{tabular}
  \qquad
  \begin{tabular}{@{}lcc@{}}
    \toprule
    Confusion & Base & guard-bear \\
    \midrule
    TN & 6   & 297 \\
    FP & 294 & 3   \\
    FN & 13  & 4   \\
    TP & 300 & 309 \\
    \bottomrule
  \end{tabular}
\end{table}

\label{sec:deployment}Table~\ref{tab:res2}'s latency figures and \textit{guard-bear}'s footprint together show a complete pipeline practical on one consumer machine; Fast-1B alone clears the 1.0s VR target, as do several cross-architecture configurations (Table~\ref{tab:crossarch}). End-to-end latency with \textit{guard-bear}'s pass included has not yet been measured (Limitations).

\subsection{Statistical Validation of the Crossover: A Multi-Seed Campaign}\label{sec:seedcampaign}

The single-run pattern above (Section~\ref{sec:responsemodel}) was an unexpected finding uncovered during development, not a designed comparison, so it needed the same reproducibility scrutiny as any surprising result before being treated as real. The ablation and the rank/layer sweeps in Section~\ref{sec:ranksweep} are single- or two-seed comparisons; Section~\ref{sec:reproducibility} documents why that scrutiny was necessary -- a 17-point swing on a nominally identical re-run, traced primarily to a dataset revision plus a smaller, unexplained non-determinism effect.

Given this, we resolved the crossover question directly with a dedicated campaign: retraining Standard and Fast at their exact defined configurations ($r=8$/\texttt{num\_layers=16} and $r=4$/\texttt{num\_layers=8}) across many seeds, one consistent (current) dataset vintage throughout, role \texttt{age\_5\_11}, FK$\leq$7.0 pass rate -- $n=25$/config at 1B (pre-committed, no interim look); at 3B, an initial $n=15$/config batch left the comparison marginal ($t=2.01$, $p=0.054$), so, per a plan fixed before unblinding stage two (a single top-up to $n=30$, not an open-ended stopping rule), we collected 15 further seeds/config and combine both stages below via a pre-specified equal-weight inverse-normal combination test, the valid correction for this design.

\begin{table}[!h]
  \centering
  \caption{Seed-campaign FK$\leq$7.0 pass rate (mean $\pm$ SD across independent training runs).}
  \label{tab:seedcampaign}
  \small
  \begin{tabular}{@{}lccc@{}}
    \toprule
    Config. & $n$ & Mean & SD \\
    \midrule
    Fast-1B     & 25 & 63.4\% & 6.3 \\
    Standard-1B & 25 & 52.1\% & 6.3 \\
    Fast-3B     & 30 & 52.2\% & 7.0 \\
    Standard-3B & 30 & 58.3\% & 6.9 \\
    \bottomrule
  \end{tabular}
\end{table}

On 1B, Fast beats Standard ($t=6.40$, $p=6.0\times10^{-8}$). On 3B, the naive pooled test gives $t=3.42$, $p=0.0011$, but the interim look makes this anti-conservative; combining that stage with the independent 15-seed replication collected afterward ($t=2.75$, $p=0.010$) via the combination test gives the valid $p=0.0015$ -- \textbf{the crossover survives the statistically correct test on both sides.} This campaign covers FK$\leq$7.0 pass rate only, for \texttt{age\_5\_11} only; the other four readability metrics, the \texttt{age\_12\_18} role, latency, and the classifier are single-run, as in Table~\ref{tab:res2}, though now on the same consistent dataset vintage as this campaign.

Because Standard and Fast differ systematically in response length (Table~\ref{tab:res2}), we tested whether the crossover survives controlling for length, across all 5{,}500 individual campaign responses. Fast is shorter than Standard at 1B (72 vs.\ 78 words, $p=4\times10^{-16}$) but longer at 3B (97 vs.\ 76 words, $p=9\times10^{-50}$). Residualizing FK grade on a single word-count regression pooled across configs (an equal-slopes assumption we did not test) and comparing residuals by config, \textbf{the 1B crossover survives length-adjustment} (residual gap 0.35 FK grade levels, $t=-5.95$, $p=3\times10^{-9}$), but \textbf{the 3B crossover is substantially a length effect} -- the residual gap shrinks to 0.10 FK grade levels, only marginal ($t=-1.92$, $p=0.055$). A joint logistic regression of pass/fail on config and word count agrees: the config coefficient stays large at 1B (0.35) but is near zero at 3B (0.007) once length is included. Standard's 3B advantage is better described as producing shorter, not more simply-worded, responses.

\subsection{Mechanism Sweeps: Rank and Layer Depth}\label{sec:ranksweep}

The Standard/Fast configurations differ in both rank and \texttt{num\_layers}, so even with the crossover confirmed (Section~\ref{sec:seedcampaign}), it is not attributable to either factor alone. The sweeps use two seeds -- lower power than the campaign above; $\pm$std is a two-point range, not a variance estimate. Both were retrained on the current vintage.

Fixing \texttt{num\_layers=8} and sweeping $r \in \{2,4,8,16\}$ across four architectures (Llama 1B/3B, Qwen~2 0.5B, SmolLM2~360M; two seeds; 32 runs; Table~\ref{tab:sweeps}, left), \textbf{only SmolLM2~360M shows a clean, monotonic rank-regime pattern} (48\%~$\to$~63\%), consistent with depth (32 layers) driving a preference for higher rank; Llama and Qwen are flat or only weakly increasing within noise. We treat SmolLM2's rank sensitivity as the sweep's one robustly replicated finding, and the broader ``rank is the operative variable'' claim as suggestive, not resolved.

A companion layer sweep (fixed \texttt{rank=4}, \texttt{num\_layers}$\in\{4,8,16\}$, Llama 1B/3B, two seeds, 12 runs; Table~\ref{tab:sweeps}, right) replicates more consistently: \textbf{3B improves with depth} (57\%~$\to$~63\%) while \textbf{1B peaks at the fewest layers} (69\% at layers=4, declining to 60\% at layers=16), matching the original sweep's qualitative conclusion. Latency scales with layer count for both sizes.

\begin{table}[!h]
  \centering
  \caption{Rank sweep (left; num\_layers=8 fixed) and layer sweep (right; rank=4 fixed) FK $\leq 7.0$ pass rate, mean $\pm$ std across two seeds (n=2; read as a range between the two runs, not a variance estimate), current dataset vintage.}
  \label{tab:sweeps}
  \scriptsize
  \setlength{\tabcolsep}{3pt}
  \begin{tabular}{@{}lcccc@{}}
    \toprule
    Model     & $r{=}2$ & $r{=}4$ & $r{=}8$ & $r{=}16$ \\
    \midrule
    Llama 1B  & 63$\pm$10 & 59$\pm$7 & \textbf{64$\pm$6} & 62$\pm$6 \\
    Llama 3B  & 54$\pm$11 & 58$\pm$3 & 57$\pm$1 & \textbf{62$\pm$11} \\
    Qwen 0.5B & 70$\pm$3 & \textbf{71$\pm$13} & 71$\pm$4 & 60$\pm$11 \\
    SmolLM2   & 48$\pm$3 & 56$\pm$6 & 62$\pm$6 & \textbf{63$\pm$4} \\
    \bottomrule
  \end{tabular}
  \qquad
  \begin{tabular}{@{}lccc@{}}
    \toprule
    Model, layers & FK$\leq$7 & Lat.(s) & Tok/s \\
    \midrule
    1B, 4  & \textbf{69\%} & 0.85 & \textbf{105.6} \\
    1B, 8  & 61\% & 0.84 & 97.4 \\
    1B, 16 & 60\% & 1.06 & 84.0 \\
    3B, 4  & 57\% & 2.73 & \textbf{45.0} \\
    3B, 8  & 60\% & 2.67 & 43.2 \\
    3B, 16 & \textbf{63\%} & \textbf{2.24} & 39.5 \\
    \bottomrule
  \end{tabular}
\end{table}

\subsection{Cross-Architecture and Capacity-Ladder Validation}\label{sec:crossarch}

To test generalization beyond Llama, we evaluate Qwen~2 0.5B and SmolLM2~360M~\cite{allal2025smollm2} and the Qwen 2.5 family (0.5B, 1.5B, 3B) for age~5--11 (Table~\ref{tab:crossarch}), single-seed and exploratory. \textbf{SmolLM2 Standard dominates} (80\% FK); \textbf{Qwen 2.5 shows Fast at or above Standard at every size}, suggesting no crossover for that family. Qwen 2 Fast and Qwen 2.5-0.5B Fast clear the 1.0s target with the most margin. The Classifier column reflects the original, pre-revision vintage.

\begin{table*}[!h]
  \centering
  \caption{Cross-architecture and Qwen 2.5 capacity-ladder results, age~5--11 (50 examples, seed 42, current dataset vintage; classifier is 5-fold CV, chance~=~0.50, $\dagger$original pre-revision vintage, not yet re-validated).}
  \label{tab:crossarch}
  \scriptsize
  \setlength{\tabcolsep}{3pt}
  \begin{tabular}{@{}lccccc@{}}
    \toprule
    Config. & FK$\leq$7.0 & Avg Lat. & $<$1.0s & Tok/s or Tok & Classifier$^\dagger$ \\
    \midrule
    Qwen 2 Std.       & 68\% & 0.54\,s & 100\% & 83 tok  & $0.940 \pm 0.049$ \\
    Qwen 2 Fast       & 62\% & 0.51\,s & 96\%  & 94 tok  & $0.960 \pm 0.058$ \\
    SmolLM2 Std.      & 80\% & 0.87\,s & 72\%  & 82 tok  & $0.950 \pm 0.032$ \\
    SmolLM2 Fast      & 60\% & 1.30\,s & 56\%  & 137 tok & $0.920 \pm 0.040$ \\
    Qwen2.5-0.5B Fast & 56\% & 0.45\,s & 100\% & 176.4   & $0.950 \pm 0.032$ \\
    Qwen2.5-0.5B Std. & 56\% & 0.57\,s & 100\% & 152.1   & $0.940 \pm 0.097$ \\
    Qwen2.5-1.5B Fast & 74\% & 1.07\,s & 58\%  & 76.1    & $\mathbf{0.980 \pm 0.024}$ \\
    Qwen2.5-1.5B Std. & 66\% & 1.28\,s & 6\%   & 67.8    & $0.890 \pm 0.073$ \\
    Qwen2.5-3B Fast   & 70\% & 1.73\,s & 4\%   & 43.0    & $0.970 \pm 0.040$ \\
    Qwen2.5-3B Std.   & 50\% & 2.18\,s & 0\%   & 39.9    & $0.930 \pm 0.068$ \\
    \bottomrule
  \end{tabular}
\end{table*}

\section{Methods}\label{sec:methodology}

\subsection{System Architecture}\label{sec:architecture}

Raw text first passes through \textit{guard-bear}, a binary input classifier; unsafe input is blocked before reaching \textit{small-bear}. Safe input is routed by an age-gate, supplied by the calling VR application, to one of two LoRA adapters on a shared frozen 4-bit quantized base model. Neither is fused into its base model. Emotional-support queries are in scope only as reassurance/redirection; clinical assessment and crisis intervention are out of scope. No system prompt is used, so the adapter weights encode register directly.

\textit{guard-bear} classifies text as \textbf{safe} (a legitimate pediatric query, including benign out-of-scope curiosity) or \textbf{unsafe} (prompt injection, jailbreak, adult/parental queries, gibberish, inappropriate content, social engineering), and is a full parameter fine-tune of a small (86M) classifier, unlike \textit{small-bear}. The design target is high recall: a missed unsafe input is worse than an incorrectly blocked legitimate query, motivating $\geq 0.97$ recall.

\subsection{Datasets}\label{subsec:responsedataset}

\textit{small-bear}'s dataset was generated by Claude Opus 4.6 in extended thinking mode from the Children's Hospital Patient and Family Guide~\cite{hospitalguide}, across five categories, each age-stratified (5--11, 12--18). An initial 1000-example release was superseded by a quality-curation pass adding a training-only \texttt{edge\_cases} category (50) and 73 diversified examples, for 1123 total, split 562/561, and 100 new validation examples; Section~\ref{sec:reproducibility} documents how this revision affected earlier results. Examples are instruction/response pairs, without a system prompt, so the model encodes register in its weights.

\textit{guard-bear}'s dataset combines synthetic safe/unsafe examples with unsafe examples from public prompt-injection and jailbreak datasets: 4{,}703 examples (2{,}404 unsafe, 2{,}299 safe) across 14 subcategories (full taxonomy in the repository). ``Safe'' means \textit{not actively harmful or adversarial}, not ``in \textit{small-bear}'s scope''; benign out-of-scope questions are handled downstream. The corpus is split 3{,}476/614/613 for training, validation, and held-out test.

\subsection{Base Models, LoRA Configuration, and Training Protocol}\label{sec:trainprotocol}

We use the Llama 3.2 Instruct family~\cite{dubey2024llama3} at 1B/3B, 4-bit quantized via \texttt{mlx\_lm}~\cite{apple2023mlx} on Apple Silicon; full-parameter fine-tuning of the 3B model exceeds the 16GB ceiling, making LoRA the only viable strategy. \textit{guard-bear} fine-tunes all 86M parameters via the Metal Performance Shaders backend.

LoRA~\cite{hu2022lora} freezes a pre-trained weight matrix $W_0 \in \mathbb{R}^{d\times k}$ and expresses updates as $\Delta W = BA$, applied to all linear projections in the targeted layers with \texttt{scale=20.0} held constant (a raw output multiplier, not an alpha/rank ratio). \textbf{Standard} ($r=8$, \texttt{num\_layers=16}) covers the upper 16 layers; \textbf{Fast} ($r=4$, \texttt{num\_layers=8}) covers the upper 8. The two differ in both rank and layer count simultaneously; Sections~\ref{sec:seedcampaign}--\ref{sec:ranksweep} confirm the resulting crossover is real.

\textit{small-bear}'s adapters train with Adam, cosine-decay LR (peak $1\times10^{-5}$, 100-step warmup, 600 steps), batch size 4, max sequence length 768, dropout 0.0, loss masked to assistant tokens; all eight runs use seed 42. \textit{guard-bear} fine-tunes \texttt{meta-llama/Prompt-Guard-86M}~\cite{promptguard2024} (DeBERTa-v2-based~\cite{he2021deberta}) on the full training split, threshold tuned post-hoc to the lowest value meeting $\geq 0.97$ recall, selecting \textbf{0.08}.

\noindent\textbf{Reproducibility and dataset vintage.}\label{sec:reproducibility} A re-run of the nominal Fast-1B configuration (identical model, data, seed) produced an FK$\leq$7.0 pass rate 17 points lower than an earlier draft's single run. A confounding library-version change was ruled out; the training and validation data had also been revised between runs (Section~\ref{subsec:responsedataset}), and re-scoring the original adapter on the revised validation set alone closed most of the gap. A smaller non-determinism effect (16--18 points, data held fixed) remains unexplained. A closer audit found the rank sweep, layer sweep, and initial cross-architecture results had also predated the revision; all were retrained on the current vintage before being reported here, except Table~\ref{tab:crossarch}'s Classifier column, still provisional. This episode motivated the dedicated multi-seed campaign in Section~\ref{sec:seedcampaign}: a single-run result cannot be trusted without independent scrutiny.

\subsection{Evaluation Framework}

For \textit{small-bear}: \textbf{readability} via \texttt{textstat}~\cite{bansal2023textstat} (Flesch-Kincaid~\cite{kincaid1975fk} primary, pass/fail at FK~$\leq$~7.0; SMOG~\cite{mclaughlin1969smog}, Gunning Fog~\cite{gunning1952fog}, Coleman-Liau~\cite{coleman1975cli}, type-token ratio); \textbf{latency}; \textbf{style separation} via TF-IDF+logistic-regression (scikit-learn~\cite{pedregosa2011sklearn}, 5-fold CV, chance~=~0.50). These are proxies, not direct measures, of age-adapted communication (Section~\ref{sec:discussion}). For \textit{guard-bear}: unsafe recall, precision, F1, accuracy, ROC-AUC, and per-subcategory recall.

\textbf{Ethical approval declarations.} Not applicable: no experiments on live vertebrates, humans, or human tissue; all data are synthetic or drawn from public reference/prompt-injection materials, with no patient contact or identifying information.

\section{Discussion}\label{sec:discussion}

For \textit{guard-bear}, the central finding is that \textbf{a safety classifier cannot be assumed to transfer to a new deployment population without evaluation}; the gap is severe, worse-than-random, not suboptimal thresholding. \textit{small-bear} without an input safety gate is not, by itself, deployable, and Section~\ref{sec:deployment} shows both are practical together.

\textbf{Fast-1B is the best Llama-family configuration} for VR deployment: the only option meeting the 1.0s latency target with strong readability and style separation; SmolLM2 Standard (Section~\ref{sec:crossarch}) beats it on FK pass rate and is a strong alternative where architecture is unconstrained. \textit{small-bear}'s most interesting outcome was not planned: the \textbf{crossover} (Section~\ref{sec:seedcampaign}) suggests small-model LoRA configuration below 7B does not follow rank-selection guidance developed at larger scales, though length-adjustment shows only the 1B side is length-independent. Its mechanism is only partially resolved (Section~\ref{sec:ranksweep}). These gains are necessary but not sufficient evidence of genuine age-adapted communication: a response can be short and lexically distinct while still emotionally mismatched to its audience, a gap only human evaluation can close.

\section{Conclusion}\label{sec:conclusion}

This paper shows two things. First, a general-purpose input-safety classifier cannot be assumed to transfer to a new, vulnerable population: \texttt{Prompt-Guard-86M}'s unmodified performance is worse than random, resolved only by dedicated fine-tuning (ROC-AUC 0.999). Second, a complete age-adaptive response and safety pipeline is achievable on one consumer machine, with measurable readability and style-separation gains over zero-shot base models (Abstract).

We also uncovered an unexpected LoRA configuration-ordering crossover between 1B and 3B Llama models (Section~\ref{sec:seedcampaign}) after an initial single-seed observation proved non-reproducible. A rank sweep and layer sweep sought to attribute this to rank and layer depth; layer-depth replicates reasonably well, but the rank-regime mechanism replicates cleanly only for SmolLM2 -- existence established, precise cause only partially so. Neither component alone is deployable: \textit{small-bear} presumes benign input, and \textit{guard-bear} makes that presumption safe.

Limitations: Section~\ref{sec:reproducibility} documents this pipeline's non-determinism and dataset-vintage confound, accounted for above; one piece remains provisional (Table~\ref{tab:crossarch}'s Classifier column). The rank sweep's mechanism claim is qualified (Section~\ref{sec:ranksweep}): SmolLM2's rank sensitivity replicates cleanly, Llama and Qwen do not, at two-seed power -- plausible, not resolved. Length-adjustment was untested on the sweeps' architectures; single-seed perplexity does not track the crossover; models above 3B remain untested. Training and validation data share one generative process, and every metric is fully automated: no human rater has judged any output. This closed evaluation loop confirms internal, not external, validity. The base-model comparison is zero-shot, the weakest baseline; neither component has clinical validation. Reported latency excludes \textit{guard-bear}'s pass; \textit{guard-bear}'s dataset has not been red-teamed or checked for leakage.

\backmatter

\bmhead{Supplementary information}

Full validation-set outputs, dataset files, and training/evaluation scripts are available at the repositories listed under Data availability below; no separate supplementary bundle accompanies this submission.

\bmhead{Acknowledgements}

Not applicable.

\section*{Declarations}

\begin{itemize}
\item \textbf{Funding:} No funding to disclose.
\item \textbf{Conflict of interest/Competing interests:} The authors declare no conflicts of interest.
\item \textbf{Ethics approval and consent to participate; consent for publication:} Not applicable; no human subjects, patient data, or clinical trial. All data are synthetic or from public reference/prompt-injection materials.
\item \textbf{Data, materials, and code availability:} \textit{small-bear}'s code/data at \url{https://github.com/abarros6/small-bear}; \textit{guard-bear}'s code/data at \url{https://github.com/abarros6/guard-bear}.
\item \textbf{Author contribution:} S.d.R.\ and R.E.\ supervised the research. A.B.\ ran all experiments, designed \textit{guard-bear}, and wrote the paper except the Introduction, written by K.K. All authors reviewed and approved the final manuscript.
\end{itemize}

\bibliography{aissh_final}

\end{document}
